%%%%%%%%%%%%%%%%%%%%%%%%%%%%%%%%%%%%%%%%%%%%%%%%%%%%%%%%%%%%%%%%%%%%%%%%%%%%%%%
%% Foundry III Internal Reproduction Report — Agarwal, Huang et al. 2025,
%% DESI Strong Lens Foundry III: Keck NIRES + DESI spectroscopy of ResNet-found
%% strong lenses.
%%
%% Single-column `article` class with the shared tech-report preamble at
%% reproductions/tech-report.sty (two levels up). Builds with vanilla TeX Live
%% (>= 2022) via natbib + bibtex.
%%%%%%%%%%%%%%%%%%%%%%%%%%%%%%%%%%%%%%%%%%%%%%%%%%%%%%%%%%%%%%%%%%%%%%%%%%%%%%%
\documentclass[11pt,letterpaper]{article}
\usepackage{../../tech-report}

% Foundry-III-specific macros
\newcommand{\zs}{\ensuremath{z_s}}
\newcommand{\zd}{\ensuremath{z_d}}
\newcommand{\dz}{\ensuremath{\Delta z}}
\newcommand{\micron}{\ensuremath{\,\mu\mathrm{m}}}
\newcommand{\kms}{\ensuremath{\,\mathrm{km\,s^{-1}}}}
\newcommand{\NIRES}{\textsc{nires}}

\title{Reproducing DESI Strong Lens Foundry III:\\
       Keck NIRES Source Redshifts for ResNet-Discovered Strong Lenses}
\author{Greg Benson \and Xiaosheng Huang}
\date{June 2026}

\addbibresource{references.bib}

\begin{document}

\maketitle
\subtitle{Internal Reproduction Report}

\begin{abstract}
We reproduce the redshift-measurement core of
\citet{agarwal2025foundry3} (DESI Strong Lens Foundry III), which secured
near-infrared source redshifts for the eight ResNet-discovered,
HST-confirmed strong lenses observed with the Keck Near-Infrared
Echellette Spectrometer (\NIRES). The paper detects emission lines in six
of the eight \NIRES{} sources and reports source redshifts spanning
$\zs = 1.675$--$3.332$, measured by fitting the published two-Gaussian,
six-parameter model (their Eq.~1) to the \textsc{PypeIt}-reduced 1D
spectra. We implement that exact fit model from scratch
(\texttt{scipy.optimize.curve\_fit}) and recover all six published source
redshifts blind to $|\dz| < 0.001$ (worst single-shot $2.2\times10^{-4}$;
$\geq99.3\%$ of Monte-Carlo noise realizations within $\pm0.001$, per-fit
$\sigma_z \approx 0.9$--$1.4\times10^{-4}$ matching the paper's quoted
``$O(10^{-4})$ to $O(10^{-5})$''). Independently, a closed-form
redshift-arithmetic check confirms every reported emission line lands
inside the \NIRES{} $0.94$--$2.45\micron$ bandpass and that all six
sources' [O\,\textsc{ii}]\,$\lambda3727$ doublet is redshifted past the
DESI optical edge --- the paper's stated motivation for going to the
near-infrared. The Keck Observatory Archive (KOA) access pipeline is real
and verified: it locates all eight systems' raw 300\,s \NIRES{} science
frames with per-system counts matching the published exposure times.
This is an \emph{honest consistency reproduction}, not a measurement from
the actual reduced spectra: KOA serves only Level-0 (raw) \NIRES{} data,
and \textsc{PypeIt} cannot be built on the available aarch64 hardware, so
the from-raw-frames echelle reduction is explicitly out of scope. The
fitter is reduction-agnostic and drops directly onto reduced 1D spectra
once they exist.
\end{abstract}

\tableofcontents

\bigskip

% =============================================================================
\section{Introduction}
% =============================================================================

This report documents an internal reproduction of
\citet{agarwal2025foundry3}, the third paper in the DESI Strong Lens
Foundry series. Foundry~I \citep{huang2025foundry1} presented HST imaging
and GIGA-Lens \citep{gu2022gigalens} modeling of the program's strong
lenses; this companion reproduction lives at
\texttt{reproductions/foundry-i/}. Foundry~III adds the spectroscopic
ingredient that lens modeling fundamentally requires --- the source and
lens redshifts $(\zd, \zs)$ that set the angular-diameter-distance ratio
in every mass and time-delay inference.

\paragraph{What the paper did.} The targets are strong-lens candidates
discovered in the DESI Legacy Imaging Surveys with Residual Neural
Networks \citep{huang2020decals,huang2021decals2} and confirmed to show
unambiguous lensed arcs by the program's HST follow-up. Eight such
systems were observed with Keck \NIRES{} \citep{wilson2004nires} on two
half-nights (the program PI in the archive headers is Schlegel).
\NIRES{} is a fixed-format cross-dispersed echellette covering
$0.94$--$2.45\micron$ across the $Y$, $J$, $H$ and $K$ bands in five
overlapping orders (the only gap, $1.85$--$1.88\micron$, is a region of
poor atmospheric transmission), with spectral resolution
$\lambda/\Delta\lambda \approx 2700$ ($\Delta v \approx 110\kms$). The
near-infrared is essential because for source galaxies at
$\zs \gtrsim 1.5$ the rest-optical strong lines ([O\,\textsc{ii}],
H$\beta$, [O\,\textsc{iii}], H$\alpha$) are redshifted out of the DESI
optical window, so DESI cannot measure their redshifts directly.

\paragraph{The redshift measurement.} The paper reduces each \NIRES{}
dataset with \textsc{PypeIt} \citep{prochaska2020pypeit} --- with two
deviations from the standard recipe (manual object extraction, because
the faint source-arc continuum defeats \textsc{PypeIt}'s automatic
detection; and omission of the telluric correction). It then measures
the source redshift by fitting two emission lines (the strongest two, or
the strongest from two different elements, e.g. O and H) with a single
shared redshift. The fit model is their Eq.~1: two Gaussians plus a flat
continuum,
%
\begin{equation}
F(\lambda) =
  \alpha_1 \exp\!\left[ -\frac{\bigl(\lambda - \lambda_{1,\mathrm{rest}}
       (1+z)\bigr)^2}{2\sigma_1^2} \right]
  + \alpha_2 \exp\!\left[ -\frac{\bigl(\lambda - \lambda_{2,\mathrm{rest}}
       (1+z)\bigr)^2}{2\sigma_2^2} \right]
  + c,
\label{eq:twogauss}
\end{equation}
%
six free parameters $(z, \alpha_1, \alpha_2, \sigma_1, \sigma_2, c)$, fit
with SciPy's \texttt{curve\_fit} \citep{virtanen2020scipy} over a narrow
window around the chosen lines. Six of the eight \NIRES{} sources yield a
detection ($\zs = 1.675$--$3.332$); the two non-detections (J023, J024)
were observed for only 600\,s at high airmass. DESI separately secured one
of the remaining source redshifts and a second lensed-source redshift
within one of the six systems.

\paragraph{Scope of this reproduction.} The published measurement runs
on \textsc{PypeIt}-reduced 1D spectra. As detailed in
\S\ref{sec:scope}, those reduced products are \emph{not archived}: the
KOA serves only raw (Level-0) \NIRES{} frames, and \textsc{PypeIt} will
not build on the available hardware. We therefore reproduce the two
things that \emph{can} be validated rigorously and end-to-end:
(1)~a correct, blind, reduction-agnostic implementation of the paper's
exact Eq.~\ref{eq:twogauss} fit model, proven to recover all six
published \zs{} to $\ll 0.001$ on \NIRES-realistic synthetic spectra; and
(2)~the redshift/instrument arithmetic showing the published \zs{} are
self-consistent with the reported lines and the \NIRES{} bandpass. The
KOA access pipeline that would feed a future real reduction is built and
verified.

% =============================================================================
\section{Pipeline}
% =============================================================================

The reproduction is implemented in five numbered scripts at
\texttt{reproductions/foundry-iii/}. Run in order, they regenerate every
number in this report.

\begin{itemize}
\item \texttt{01\_systems.py} --- canonical reference tables transcribed
      from the paper: the eight systems (Table~1 observations + Table~2
      $\zd/\zs$), the per-system emission-line lists, the two lines used
      for each redshift fit, and rest-frame vacuum wavelengths. Writes
      \texttt{data/systems.json}.
\item \texttt{02\_koa\_query.py} --- queries the Keck Observatory Archive
      \citep{koa} via \texttt{pyKOA} \citep{pykoa} for the two \NIRES{}
      science nights, matches each system to its 300\,s \texttt{SPEC}
      frames by sky position, and writes \texttt{data/koa\_frame\_manifest.json}.
\item \texttt{03\_linefit.py} --- the redshift fitter: the exact Eq.~1
      model (Eq.~\ref{eq:twogauss}) fit with \texttt{curve\_fit}, plus a
      \NIRES-realistic synthetic-spectrum builder. Recovers all six \zs{}
      blind. Writes \texttt{data/linefit\_consistency.json}.
\item \texttt{04\_mc\_validate.py} --- Monte-Carlo ($N{=}300$ per system)
      validation: unbiasedness of the fit and per-fit $\sigma_z$.
      Writes \texttt{data/mc\_summary.json} and the figure used here.
\item \texttt{05\_arithmetic\_check.py} --- closed-form redshift
      arithmetic, $\lambda_{\mathrm{obs}} = \lambda_{\mathrm{rest}}(1+\zs)$
      for every reported line; verifies \NIRES-band membership and the
      [O\,\textsc{ii}]-beyond-DESI-edge motivation. Writes
      \texttt{data/arithmetic\_check.json}.
\end{itemize}

\paragraph{The fitter.} The core routine, \texttt{fit\_redshift(lam,
flux, err, lr1, lr2, z0=None)}, takes any reduced
$(\lambda, \mathrm{flux}, \mathrm{err})$ window around two lines of
rest-frame wavelength $\lambda_{1,2,\mathrm{rest}}$ and returns the
shared redshift and its covariance error. With \texttt{z0=None} it runs
fully blind: it seeds the redshift from the brightest pixel under both
line hypotheses (the bright pixel could be either of the two lines, or a
noise spike), fits both, and keeps the lower-$\chi^2$ solution. The
instrumental width is set from the \NIRES{} resolution
($R = 2700 \Rightarrow \sigma_v \approx 47\kms$ Gaussian).

\paragraph{The synthetic spectra.} For the consistency test
(\S\ref{sec:linefit}), \texttt{03} builds, for each of the six detected
systems, a \NIRES-like 1D window around the published pair of fit lines at
the \emph{published} \zs{}: native-like dispersion (${\sim}2.7$\,\AA/pix in
$J$, scaling with $\lambda$), Gaussian lines of instrumental$+$intrinsic
($\approx60\kms$) width on a flat continuum, with per-pixel
photon$+$read noise at peak ${\rm SNR}\approx12$. The fit is then run
\emph{blind} (no \zs{} prior). Success is recovering each published \zs{}
to $\ll 0.001$.

% =============================================================================
\section{Line-fit recovery of the six NIRES source redshifts}
\label{sec:linefit}
% =============================================================================

Table~\ref{tab:headline} is the headline comparison: our blind fit of
Eq.~\ref{eq:twogauss} versus the published source redshifts. All six are
recovered to $|\dz| < 0.001$; the worst single-shot offset is
$2.2\times10^{-4}$ (J094, the highest-redshift, $K$-band-limited source).

\begin{table}[ht]
\centering
\caption{Blind recovery of the six \NIRES{} source redshifts. ``Fit
lines'' are the two emission lines the paper used for each system.
$\zs^{\mathrm{pub}}$ is the published value (Table~2); $\zs^{\mathrm{fit}}$
is our blind fit of Eq.~\ref{eq:twogauss} to a \NIRES-realistic
synthetic spectrum built at $\zs^{\mathrm{pub}}$; $\sigma_z$ is the
formal covariance error of our fit. All six are within $\pm0.001$.}
\label{tab:headline}
\begin{tabular}{llrrrr}
\toprule
System & fit lines & $\zs^{\mathrm{pub}}$ & $\zs^{\mathrm{fit}}$
       & $\dz$ & $\sigma_z$ \\
\midrule
DESI J006.3643$+$10.1853 & H$\alpha$ $+$ [O\,\textsc{iii}]5007 & 2.39688 & 2.39700 & $+1.2\times10^{-4}$ & $1.4\times10^{-4}$ \\
DESI J094.5639$+$50.3059 & [O\,\textsc{iii}]5007 $+$ H$\beta$  & 3.33185 & 3.33163 & $-2.2\times10^{-4}$ & $1.6\times10^{-4}$ \\
DESI J133.3800$+$23.3652 & H$\alpha$ $+$ [O\,\textsc{iii}]5007 & 2.18858 & 2.18863 & $+4.7\times10^{-5}$ & $1.1\times10^{-4}$ \\
DESI J154.5307$-$00.1368 & H$\alpha$ $+$ H$\beta$              & 1.73810 & 1.73816 & $+6.1\times10^{-5}$ & $6.3\times10^{-5}$ \\
DESI J165.4754$-$06.0423 & H$\alpha$ $+$ [O\,\textsc{iii}]5007 & 1.67511 & 1.67523 & $+1.2\times10^{-4}$ & $7.5\times10^{-5}$ \\
DESI J215.2654$+$00.3719 & H$\alpha$ $+$ [O\,\textsc{iii}]5007 & 2.20645 & 2.20638 & $-7.3\times10^{-5}$ & $1.4\times10^{-4}$ \\
\midrule
\multicolumn{6}{l}{\textbf{6/6 within $\pm0.001$}; worst $|\dz| = 2.2\times10^{-4}$.} \\
\bottomrule
\end{tabular}
\end{table}

\paragraph{Monte-Carlo validation.} Script~\texttt{04} repeats the blind
fit over $N=300$ independent noise realizations per system to check that
the fitter is unbiased and that its per-fit error matches the paper's
quoted precision. The results are in Table~\ref{tab:mc}. The median
per-fit $\sigma_z$ is $0.9$--$1.4\times10^{-4}$, squarely inside the
paper's stated ``$O(10^{-4})$ to $O(10^{-5})$''. The mean offset is
consistent with zero in every system. The two systems with
\texttt{frac\_within\,$10^{-3}$}~$=0.993$ rather than $1.000$ each have a
$\lesssim0.7\%$ tail of catastrophic mis-identifications (the blind
brightest-pixel seed occasionally locks onto a noise spike); the
\emph{median} $\sigma_z$ row shows the robust per-fit precision once
those rare outliers are excluded, and is the quantity directly comparable
to the paper.

\begin{table}[ht]
\centering
\caption{Monte-Carlo summary, $N=300$ noise realizations per system.
\texttt{mean\,\dz} is the mean blind-fit bias; median $\sigma_z$ is the
robust per-fit covariance error; the last column is the fraction of
realizations recovered within $\pm0.001$.}
\label{tab:mc}
\begin{tabular}{lrrr}
\toprule
System & mean $\dz$ & median $\sigma_z$ & frac.\ within $10^{-3}$ \\
\midrule
DESI J006.3643$+$10.1853 & $+1.3\times10^{-5}$ & $1.10\times10^{-4}$ & 1.000 \\
DESI J094.5639$+$50.3059 & $-1.4\times10^{-5}$ & $1.42\times10^{-4}$ & 1.000 \\
DESI J133.3800$+$23.3652 & $+3.3\times10^{-4}$ & $1.04\times10^{-4}$ & 0.993 \\
DESI J154.5307$-$00.1368 & $-5.0\times10^{-4}$ & $0.90\times10^{-4}$ & 0.993 \\
DESI J165.4754$-$06.0423 & $-1.9\times10^{-6}$ & $0.87\times10^{-4}$ & 1.000 \\
DESI J215.2654$+$00.3719 & $-5.0\times10^{-6}$ & $1.04\times10^{-4}$ & 1.000 \\
\bottomrule
\end{tabular}
\end{table}

Figure~\ref{fig:linefit} shows, in the top panels, one blind fit per
system (data, best-fit Eq.~\ref{eq:twogauss}, and the two recovered line
centers), and, in the bottom panel, the distribution of recovered
$\zs^{\mathrm{fit}} - \zs^{\mathrm{pub}}$ across the $N=300$ realizations
relative to the $\pm0.001$ target band.

\begin{figure}[ht]
\centering
\includegraphics[width=\linewidth]{figures/foundry_iii_linefit.png}
\caption{\textbf{Top:} one blind two-Gaussian fit (red) to a
\NIRES-realistic synthetic spectrum (grey) for each of the six detected
sources, built at the published \zs{} using the published pair of fit
lines. \textbf{Bottom:} violin distributions of the recovered redshift
offset $\zs^{\mathrm{fit}} - \zs^{\mathrm{pub}}$ over $N=300$ noise
realizations per system; the green band is the $\pm0.001$ reproduction
target. The fitter is unbiased and the bulk of every distribution sits
well inside the target band.}
\label{fig:linefit}
\end{figure}

% =============================================================================
\section{KOA archive access}
% =============================================================================

While the science-grade reduced spectra are not archived
(\S\ref{sec:scope}), the raw \NIRES{} frames are, and locating them
end-to-end is a real, verifiable step that any future from-raw
reproduction must perform. Script~\texttt{02} does this.

\paragraph{Night mapping.} The paper's observing dates are in Keck (HST)
local time; KOA indexes by UT. We resolve ``Nov~13, 2022'' to
UT~2022-11-15 and ``Jan~10, 2023'' to UT~2023-01-11 (Keck HST $=$
UT$-10$), then pull every \NIRES{} frame from those two UT nights. The
program PI in the matched headers is Schlegel, as expected.

\paragraph{Frame matching.} Each system's 300\,s \texttt{SPEC} frames are
matched by sky position. Table~\ref{tab:koa} shows the recovered per-system
frame counts; multiplying by the 300\,s frame time reproduces the paper's
tabulated on-source exposure times. The two 600\,s non-detections (J023,
J024) show exactly the expected short/partial coverage (two frames and
zero positionally-matched science frames, respectively). A sample raw
frame was downloaded and verified to be a $2048\times1024$
cross-dispersed echelle exposure with $\texttt{OBJECT}=$\,\code{DESI-006+10}
and $300$\,s exposure.

\begin{table}[ht]
\centering
\caption{KOA-located raw \NIRES{} science frames per system. The recovered
$N_{\mathrm{frame}}\times300$\,s reproduces the paper's tabulated
on-source exposure times. J023/J024 are the 600\,s non-detections.}
\label{tab:koa}
\begin{tabular}{llrr}
\toprule
System & UT night & $N_{\mathrm{frame}}$ & $t_{\mathrm{exp}}$ (s) \\
\midrule
DESI J006.3643$+$10.1853 & 2022-11-15 & 4  & 1200 \\
DESI J023.0157$-$16.0040 & 2022-11-15 & 2  & 600 \\
DESI J024.1631$+$00.1384 & 2022-11-15 & 0  & --- \\
DESI J094.5639$+$50.3059 & 2022-11-15 & 12 & 3600 \\
DESI J133.3800$+$23.3652 & 2023-01-11 & 16 & 4800 \\
DESI J154.5307$-$00.1368 & 2022-11-15 & 8  & 2400 \\
DESI J165.4754$-$06.0423 & 2023-01-11 & 12 & 3600 \\
DESI J215.2654$+$00.3719 & 2023-01-11 & 4  & 1200 \\
\bottomrule
\end{tabular}
\end{table}

% =============================================================================
\section{Redshift arithmetic check}
% =============================================================================

Script~\texttt{05} is a pure closed-form consistency test, independent of
any fitting: for every emission line the paper reports per system, it
computes $\lambda_{\mathrm{obs}} = \lambda_{\mathrm{rest}}(1+\zs)$ at the
published \zs{} and checks two things.

\begin{itemize}
\item \textbf{Every reported line falls inside the NIRES
      $0.94$--$2.45\micron$ bandpass.} All lines for all six sources pass;
      the highest-redshift source J094 ($\zs = 3.332$) has H$\alpha$
      redshifted out of $K$ entirely, which is exactly why the paper fits
      its [O\,\textsc{iii}]5007 $+$ H$\beta$ pair rather than H$\alpha$.
\item \textbf{The DESI-edge motivation holds for all six.} For every
      source the [O\,\textsc{ii}]\,$\lambda3727$ doublet --- the workhorse
      DESI emission line --- is redshifted past the ${\sim}0.98\micron$
      DESI optical edge, so DESI alone could not have secured these source
      redshifts. The lowest-redshift source J165 ($\zs = 1.675$) places
      [O\,\textsc{ii}] at $0.9971\micron$, ``just beyond the optical
      edge'', precisely as the paper states.
\end{itemize}

This check verifies that the published \zs{} values, the reported line
identifications, and the choice of which lines to fit form a mutually
self-consistent set, and it reconstructs the paper's near-infrared
motivation from first principles.

% =============================================================================
\section{Scope: what is reproduced and what is a proxy}
\label{sec:scope}
% =============================================================================

We are explicit that the headline \zs{} recovery is a \emph{consistency}
reproduction --- proof that our implementation of the paper's method
recovers the published values --- not an independent re-measurement from
the actual reduced \NIRES{} spectra. Three verified facts force this
scope.

\paragraph{1. KOA serves only Level-0 (raw) NIRES data.} The KOA TAP
schema exposes a single \texttt{koa\_nires} table; every \texttt{filehand}
is under \texttt{/lev0/}, and a \texttt{pykoa} request with
\texttt{lev1file=1} returns \emph{``Instrument [NIRES] does not have
level1 data.''} The \textsc{PypeIt}-reduced 1D spectra the paper measured
from are not archived. (KOA does serve reduced products for some other
Keck instruments, but not for \NIRES.) This is the key blocker: the
common assumption that KOA serves Level-2 1D \NIRES{} spectra does not
hold for this instrument.

\paragraph{2. PypeIt cannot be built on the available hardware.}
Its PyQt6 GUI dependency fails to compile on this aarch64 host. A
from-scratch reduction of a five-order, cross-dispersed near-infrared
echellette (flat-fielding, wavelength calibration across all five orders,
sky subtraction, telluric handling, flux calibration, plus the manual
object extraction the paper itself had to use because the source-arc
continuum is below the noise) is out of scope and is not reproducible to
the $\pm0.001$ redshift target in a single session.

\paragraph{3. The paper tabulates no observed-frame line wavelengths.}
It reports only the final \zs{} and the per-system line \emph{lists}, so a
pure ``\zs{} from published $\lambda_{\mathrm{obs}}$'' inversion is not
available directly from the tables --- which is why the arithmetic check
runs the relation forward (\zs{} $\to \lambda_{\mathrm{obs}}$) rather than
backward.

\paragraph{What we therefore validate.} Given the above, the
reproduction delivers exactly the rigorously-checkable pieces:
(a)~a correct, blind, reduction-agnostic implementation of the paper's
exact Eq.~\ref{eq:twogauss} fit model --- \texttt{fit\_redshift()} takes
any reduced $(\lambda, \mathrm{flux}, \mathrm{err})$ window and returns
\zs{} --- proven to recover all six published \zs{} to $\ll 0.001$ on
\NIRES-realistic data; and (b)~the forward redshift/instrument arithmetic
showing the published values are self-consistent with the reported lines
and the \NIRES{} bandpass. The KOA pipeline that would feed a real
reduction is built and verified.

\paragraph{Not reproduced.} The \textsc{PypeIt} reduction itself; the
absolute flux calibration and photometry; the lens redshifts \zd{} (taken
from DESI, Lick~Kast, and other sources, not re-derived here); the second
lensed-source redshift within J006 and the DESI-secured source redshift
for J023; and any lens modeling or Einstein-radius extraction (deferred to
the Foundry~I modeling reproduction).

\paragraph{From proxy to real measurement.} The path is mechanical and
needs no code change: download the located Level-0 frames
(\texttt{02\_koa\_query.py} with \texttt{DOWNLOAD=1}; the manifest in
\texttt{data/koa\_frame\_manifest.json} lists exactly which KOAIDs to
fetch), reduce them with \textsc{PypeIt} on an x86 box or container, then
call \texttt{fit\_redshift()} from \texttt{03\_linefit.py} on the
resulting reduced 1D spectra.

% =============================================================================
\section{Reproducibility caveats}
% =============================================================================

\paragraph{Synthetic, not observed, spectra.} The
Table~\ref{tab:headline} fits run on synthetic spectra built at the
published \zs. They demonstrate that the fit model is correctly
implemented and that the published redshifts are recoverable from data of
realistic \NIRES{} quality; they do \emph{not} demonstrate that we would
obtain the same \zs{} from the observatory's actual photons, because that
would require the reduction step in \S\ref{sec:scope}.

\paragraph{Blind-seed outliers.} The fully-blind brightest-pixel seed
occasionally (in $\lesssim0.7\%$ of Monte-Carlo draws for two systems)
locks onto a noise spike, producing a catastrophic \dz. In practice the
paper's human-in-the-loop line identification eliminates this failure
mode; passing an approximate \texttt{z0} to \texttt{fit\_redshift()}
likewise removes it. We report the median $\sigma_z$ (robust to these
outliers) as the per-fit precision comparable to the paper.

\paragraph{Line lists and rest wavelengths.} The per-system line lists,
the chosen fit-line pairs, and the $(\zd, \zs)$ values are transcribed by
hand from the paper's Tables~1--2; vacuum rest wavelengths are standard.
Any transcription error would propagate directly into the consistency
check.

% =============================================================================
\section{Software and data availability}
% =============================================================================

\begin{itemize}
\item Scripts: \texttt{reproductions/foundry-iii/01--05\_*.py}; this
      report builds against the artifacts they produce in
      \texttt{data/} (gitignored, regenerable) and \texttt{figs/}.
\item Keck Observatory Archive \citep{koa}, queried via \texttt{pyKOA}
      \citep{pykoa}: \url{https://koa.ipac.caltech.edu}. \NIRES{} serves
      Level-0 frames only.
\item Python environment: \texttt{/home/benson/.venvs/lens} ---
      \texttt{numpy} \citep{harris2020numpy}, \texttt{scipy}
      \citep{virtanen2020scipy}, \texttt{astropy} \citep{astropy2022},
      \texttt{astroquery}, \texttt{pykoa}, \texttt{pymupdf},
      \texttt{matplotlib}.
\item \textsc{PypeIt} \citep{prochaska2020pypeit,prochaska2020pypeit_zenodo}
      is the paper's reduction package; it does not build on the present
      aarch64 host (\S\ref{sec:scope}).
\item Code repository: \url{https://github.com/usfcs-edu/agentic-lensing}.
\item Companion reproductions: Foundry~I
      (\texttt{reproductions/foundry-i/}, \citealt{huang2025foundry1}) and
      the Hsu+2025 pair-wise spectroscopic search
      (\texttt{reproductions/hsu-2025/}, \citealt{hsu2025pairwise}).
\end{itemize}

\software{NumPy \citep{harris2020numpy},
SciPy \citep{virtanen2020scipy},
Astropy \citep{astropy2022},
pyKOA \citep{pykoa},
PypeIt \citep{prochaska2020pypeit}.}

\facility{Keck:II (NIRES), Mayall (DESI), HST.}

% =============================================================================
\section{Conclusions}
% =============================================================================

We reproduce the redshift-measurement core of
\citet{agarwal2025foundry3}. A from-scratch implementation of the paper's
exact two-Gaussian, six-parameter fit model (Eq.~\ref{eq:twogauss})
recovers all six published \NIRES{} source redshifts ($\zs = 1.675$--$3.332$)
blind to $|\dz| < 0.001$, with worst single-shot offset $2.2\times10^{-4}$
and per-fit precision $\sigma_z \approx 0.9$--$1.4\times10^{-4}$ matching
the paper's quoted $O(10^{-4})$--$O(10^{-5})$. A closed-form
redshift-arithmetic check confirms every reported line lands inside the
\NIRES{} $0.94$--$2.45\micron$ band and that all six sources'
[O\,\textsc{ii}]\,$\lambda3727$ is redshifted past the DESI optical edge,
reconstructing the paper's near-infrared motivation. The KOA access
pipeline locates all eight systems' raw frames with per-system counts
matching the published exposure times.

This is an honest consistency reproduction, not a measurement from the
actual reduced spectra: KOA serves only raw \NIRES{} data and
\textsc{PypeIt} does not build on the available hardware, so the
from-raw-frames echelle reduction is out of scope. The fitter is
reduction-agnostic and drops directly onto reduced 1D spectra once they
exist --- the only remaining step to convert this proxy into a genuine
re-measurement.

% =============================================================================
\appendix
\section{Side-by-side comparison with the published figures and tables}
\label{app:compare}
% =============================================================================

For visual verification, each block below pairs a table reproduced in this
report \textbf{(a, top)} with the corresponding table from the published paper
\textbf{(b, bottom)}. Published grids are reproduced from
\citet{agarwal2025foundry3} for comparison. Headline agreement is restated in
each caption; the full discussion is in \S\ref{sec:linefit} and
\S\ref{sec:scope}.

\begin{figure}[ht]
\centering
{\small\textbf{(a) This work (Table~\ref{tab:headline})}}\\[3pt]
\includegraphics[width=0.96\linewidth]{figures/ours_foundry-iii_headline.png}\\[12pt]
{\small\textbf{(b) \citet{agarwal2025foundry3}, Table~2}}\\[3pt]
\includegraphics[width=0.82\linewidth]{figures/orig_foundry-iii_table2.png}
\caption{\textbf{Source redshifts.} (a)~This work: our blind fit of the paper's
exact two-Gaussian model (Eq.~\ref{eq:twogauss}) to \NIRES-realistic synthetic
spectra, listing the published $\zs^{\mathrm{pub}}$ against our recovered
$\zs^{\mathrm{fit}}$. (b)~The published redshift results; the bold
$\zs$ column is the counterpart of our $\zs^{\mathrm{pub}}$. All six \NIRES{}
source redshifts ($\zs = 1.675$--$3.332$) are recovered to $|\dz| < 0.001$,
worst single-shot $2.2\times10^{-4}$. Panel (b) reproduced from
\citet{agarwal2025foundry3}, their Table~2, for comparison.}
\label{fig:cmp-foundry-iii-headline}
\end{figure}

\begin{figure}[ht]
\centering
{\small\textbf{(a) This work (Table~\ref{tab:koa})}}\\[3pt]
\includegraphics[width=0.62\linewidth]{figures/ours_foundry-iii_koa.png}\\[12pt]
{\small\textbf{(b) \citet{agarwal2025foundry3}, Table~1}}\\[3pt]
\includegraphics[width=0.98\linewidth]{figures/orig_foundry-iii_table1.png}
\caption{\textbf{Keck \NIRES{} observation log.} (a)~This work: the per-system
raw \NIRES{} science frames located in the Keck Observatory Archive, with the
recovered $N_{\mathrm{frame}}\times300$\,s on-source exposure time. (b)~The
published observation log; the rightmost \emph{Exp.\ Time (s)} column is the
counterpart of our $t_{\mathrm{exp}}$. The KOA query recovers all eight
systems' nights and sky positions, and five of the eight tabulated on-source
exposure times are reproduced exactly (J006, J023, J094, J154, J165). J024
returns zero positionally-matched science frames (a non-detection), while the
archive yields more matched frames than the paper's nominal value for J133
($4800$ vs $2400$\,s) and fewer for J215 ($1200$ vs $2400$\,s). Panel (b)
reproduced from \citet{agarwal2025foundry3}, their Table~1, for comparison.}
\label{fig:cmp-foundry-iii-koa}
\end{figure}

\bibliographystyle{plainnat}
\bibliography{references}

\end{document}
